% Generated draft: requires mathematical and source review.
\documentclass[11pt]{article}
\usepackage[margin=25mm]{geometry}
\usepackage{fontspec}
\setmainfont{Latin Modern Roman}
\usepackage{amsmath,amssymb,mathtools,unicode-math}
\setmathfont{Latin Modern Math}
\usepackage{enumitem,needspace,xurl,hyperref}
\setlength{\parindent}{0pt}
\setlength{\parskip}{5pt}
\setlength{\emergencystretch}{3em}
\newcommand{\N}{\mathbb{N}}
\newcommand{\Z}{\mathbb{Z}}
\newcommand{\Q}{\mathbb{Q}}
\newcommand{\R}{\mathbb{R}}
\newcommand{\C}{\mathbb{C}}
\newcommand{\F}{\mathbb{F}}
\newcommand{\A}{\mathbb{A}}
\newcommand{\PP}{\mathbb{P}}
\newcommand{\E}{\mathbb{E}}
\newcommand{\eps}{\varepsilon}
\newcommand{\dd}{\,\mathrm{d}}
\begin{document}
\section*{UnsolvedMath ID: 30002832}
\textbf{Source problem code:} catalogue entry
\subsection{Definitions and mathematical statement}
A \emph{Fano polygon} is a convex lattice polygon \(P\subset N_{\R}\), where \(N\cong\Z^2\), whose vertices are primitive lattice points and whose strict interior contains the origin. Write \(X_P\) for the toric surface whose fan has rays through the vertices of \(P\) and two-dimensional cones over its edges. A \(\Q\)-Gorenstein deformation, abbreviated qG-deformation, is a deformation for which the relative canonical divisor is \(\Q\)-Cartier. The base field and the precise deformation-family conventions are not specified in the supplied record.

Retain the source notation \(L_P^T\) for its Laurent-polynomial parameter space, \(\pi_P\) for its classical-period map, \(H_X^{\mathrm{ts}}\) for its quantum-period parameter space, and \(\widehat G_X\) for its regularized quantum-period map. The record does not define these spaces, their coefficient fields, or their normalization conventions. Likewise, \(\operatorname{SC}(P)\) below denotes the source's singularity content. Polygon mutation, class-TG del Pezzo surfaces, and qG-deformation equivalence are understood in the cited source's conventions; their full definitions are absent from the supplied excerpt. In particular, the superscripts \(T\) and \(\mathrm{ts}\) must not be assigned additional conditions without consulting that source. An affine-linear map has the form \(v\mapsto Av+b\); an affine-linear identification means an affine-linear isomorphism.

The supplied record attributes the following three conjectural assertions to the period and deformation conjectures B--D of Akhtar--Coates--Galkin--Kasprzyk.
\begin{enumerate}
\item[(B)] For every Fano polygon \(P\) and every \(X\) satisfying the source's qualification that it is a generic qG-deformation of \(X_P\), there exists an affine-linear map
\[
\phi\colon L_P^T\longrightarrow H_X^{\mathrm{ts}}
\]
such that
\[
\widehat G_X\circ\phi=\pi_P.
\]
Thus the asserted identity is an identity of period maps throughout \(L_P^T\), not merely equality at a selected parameter value. The record does not specify the deformation space in which \(X\) is generic or the meaning of that genericity condition, and does not explicitly require \(\phi\) to be an isomorphism.

\item[(C)] For every pair of Fano polygons \(P_1,P_2\) with
\[
\operatorname{SC}(P_1)=\operatorname{SC}(P_2),
\]
if there exists an affine-linear identification
\[
\alpha\colon L_{P_1}^T\xrightarrow{\sim}L_{P_2}^T
\quad\text{such that}\quad
\pi_{P_2}\circ\alpha=\pi_{P_1},
\]
then \(P_1\) and \(P_2\) are mutation equivalent: \(P_2\) can be obtained from \(P_1\) by a finite sequence of the polygon mutations allowed in the source. The treatment of lattice isomorphisms within this equivalence is not specified in the record.

\item[(D)] For every pair of class-TG del Pezzo surfaces \(X_1,X_2\) having the same qG-rigid cyclic quotient singularities, equality of their quantum periods implies that \(X_1\) and \(X_2\) are qG-deformation equivalent. Here qG-rigid means admitting no nontrivial qG-deformations. The record does not specify whether the period comparison concerns specialized series or parameterized period maps, nor how any parameters are identified. Consequently, this last implication cannot be made into an unqualified equality of explicitly defined maps from the supplied information alone. The meanings of class-TG, sameness of the singularities, and the permitted qG-deformation families must also be recovered from the cited source.
\end{enumerate}
All three assertions retain the source's hypotheses and period conventions. The missing definitions and comparison conventions are unresolved qualifications of this statement-only draft.
\subsection{Short English statement}
The three conjectures ask whether affine-linear mirror maps identify classical and quantum periods, whether matching classical-period families determine a Fano polygon up to mutation when singularity content is fixed, and whether matching quantum periods determine a class-TG del Pezzo surface up to \(\Q\)-Gorenstein deformation when the specified rigid singularities are fixed. Several essential source definitions and period-comparison conventions are missing, so this draft requires mathematical review.
\subsection{Sources}
\begin{itemize}
\item Supplied local problem registry: permanent UnsolvedMath ID \(30002832\), exact source problem code \texttt{catalogue entry}, title \emph{Period Conjectures for Fano Polygons and Del Pezzo Surfaces}; proposed year 2015. No individual proposer is specified.
\item \emph{Algebraic Geometry} (2015), Oberwolfach Reports, pp.\ 783--836, DOI 10.4171/owr/2015/15. Supplied report reference: \url{https://doi.org/10.4171/owr/2015/15}. The registry describes these reports as records of mathematical workshops and their open-problem discussions.
\item Akhtar--Coates--Galkin--Kasprzyk, supplied reference labeled ``conjectures B--D'': \url{https://kasprzyk.work/research/pdf/mirror_symmetry_del_pezzo.pdf}.
\item Cavey--Prince, supplied reference labeled ``verified cases'': \url{https://www.jstage.jst.go.jp/article/jmath/72/2/72_465/_pdf}.
\item Coates--Kasprzyk--Prince, supplied reference labeled ``later progress'': \url{https://pmc.ncbi.nlm.nih.gov/articles/PMC8526169/}.
\item Literature assessment dated exactly \texttt{2026-08-21}: the supplied curation labels the problem partially solved and reports that special baskets and smooth cases are established, while no general theorem proves the remaining assertions. This is supplied dated curation derived from the Oberwolfach record, not an independently verified literature assessment; the time-sensitive status and cited claims require independent verification.
\item The registry labels the statement \texttt{corrected\_verified} and reports that it was checked against the official Oberwolfach report and preserves its hypotheses, notation, and open target. These are supplied verification claims; no source was independently visited or checked in preparing this draft.
\end{itemize}
\end{document}
